% Generado por documento/generar_figuras.py. No se edita a mano.
\begin{tikzpicture}[caja/.style={draw=#1, fill=#1!7, rounded corners=2pt, align=center, font=\footnotesize,
  minimum width=3.5cm, minimum height=1.25cm, inner sep=3pt}, flecha/.style={-{Stealth[length=5pt]}, draw=neutro, thick}]
\node[caja=neutro] (steam) at (0, 0) {API de Steam\\\texttt{appreviews}\\y \texttt{appdetails}};
\node[caja=neutro] (ingesta) at (3.95, 0) {Ingesta\\\texttt{ingesta\_steam.py}};
\node[caja=neutro] (releases) at (7.9, 0) {Releases con tag fijo\\data-v1 y data-v3\\con su sha256};
\node[caja=neutro] (build) at (11.85, 0) {Build en Render\\entrena con data-v1\\y compara las bandas};
\node[caja=neutro] (notebooks) at (3.95, -1.9) {Notebooks 00, 01 y 02\\código del tag codigo-v6};
\node[caja=acento] (frontend) at (7.9, -1.9) {Frontend Angular\\en Vercel};
\node[caja=acento] (api) at (11.85, -1.9) {API FastAPI\\en Render};
\draw[flecha] (steam) -- (ingesta);
\draw[flecha] (ingesta) -- (releases);
\draw[flecha] (releases) -- (build);
\draw[flecha] (releases.south) -- (notebooks.north east);
\draw[flecha] (build) -- (api);
\draw[flecha, {Stealth[length=5pt]}-{Stealth[length=5pt]}] (frontend) -- (api);
\end{tikzpicture}
